\bookchapter{第六十四章：数据存储与网络请求}

\chapterintro{一个只能“看”的 App 是死的；能存数据、能连网络的 App 才是活的。这一章认识 iOS 数据层的“三件套”：UserDefaults（存小配置）、SwiftData（存结构化数据）、URLSession（和服务器对话），外加解析 JSON 的 Codable 和现代化的 async/await。学完这章，你就能写会记住用户习惯、会实时更新内容的真 App 了。}

\booksection{一、存储方案怎么选：先看这张表}

\begin{center}
\begin{tabularx}{\linewidth}{|X|X|X|X|}
\hline
\textbf{方案} & \textbf{适合存什么} & \textbf{特点} & \textbf{本书哪里用} \tabularnewline
\hline
UserDefaults & 零散配置（开关、偏好） & 键值对，简单直接 & 第 65 章设置、第 66 章收藏城市 \tabularnewline
SwiftData & 结构化数据（账目、记录） & 官方数据库封装，@Model 一行声明 & 第 65 章账目记录 \tabularnewline
URLSession & （不存，联网取）服务器数据 & 系统自带网络框架，无第三方依赖 & 第 66 章天气数据 \tabularnewline
\hline
\end{tabularx}
\end{center}

\begin{center}\small\color{gray}数据三件套的分工\end{center}

记一个口诀：\textbf{“小配置进 UserDefaults，成批数据进 SwiftData，外面的数据走 URLSession”}。旧教程里的 Core Data 依然能用（SwiftData 就是它的现代马甲），手写 SQLite 则基本不用碰——直接学新的，别走老路。

\booksection{二、UserDefaults：存点小配置}

\begin{lstlisting}[language=Swift]
// 写：存一个值
UserDefaults.standard.set(true, forKey: "isDarkMode")
UserDefaults.standard.set("北京", forKey: "favoriteCity")

// 读：取出来（没有就返回默认值）
let isDark = UserDefaults.standard.bool(forKey: "isDarkMode")
let city = UserDefaults.standard.string(forKey: "favoriteCity") ?? "北京"

// 删：移除一个键
UserDefaults.standard.removeObject(forKey: "isDarkMode")
\end{lstlisting}

注意它只适合\textbf{零散的小配置}——几十上百条记录请交给 SwiftData，别把 UserDefaults 当数据库用（会拖慢启动、也不便查询）。

\booksection{三、SwiftData：给数据安个“真正的家”}

SwiftData 是 iOS 17 起的官方数据库方案，把“建表、增删改查”简化到了极致——\textbf{一个 @Model 注解，数据类就变成数据库表}：

\begin{lstlisting}[language=Swift]
import SwiftData

// 1. @Model：这个类自动变成"表"
@Model
class Record {
    var amount: Double
    var category: String
    var note: String
    var date: Date

    init(amount: Double, category: String, note: String, date: Date) {
        self.amount = amount
        self.category = category
        self.note = note
        self.date = date
    }
}

// 2. App 入口挂上容器
@main
struct MyApp: App {
    var body: some Scene {
        WindowGroup {
            ContentView()
        }
        .modelContainer(for: Record.self)   // 数据库就绪
    }
}

// 3. 视图里用 @Query 读、@Environment 写
struct RecordListView: View {
    @Query(sort: \Record.date, order: .reverse) var records: [Record]
    @Environment(\.modelContext) private var context

    var body: some View {
        List(records) { r in
            Text("\(r.category) \(r.amount) 元")
        }
    }

    func add(amount: Double, category: String) {
        let r = Record(amount: amount, category: category, note: "", date: Date())
        context.insert(r)     // 存
        try? context.save()   // 落盘
    }
}
\end{lstlisting}

这套“@Model + modelContainer + @Query”模板可以原样搬进你的任何项目——换换字段名就行。数据一变，@Query 自动刷新界面，妙不可言。

\booksection{四、Codable：给 JSON“翻译”}

服务器的数据大多是 JSON 文本。Swift 的\textbf{Codable}协议让“JSON 结构体”的转换自动发生——字段名对得上，魔法就生效：

\begin{lstlisting}[language=Swift]
// JSON: {"name":"北京","main":{"temp":26.5},"weather":[{"description":"晴"}]}
struct WeatherResponse: Codable {
    let name: String
    let main: Main
    let weather: [Weather]
}
struct Main: Codable { let temp: Double }
struct Weather: Codable { let description: String }

// JSON 字符串  结构体
let decoder = JSONDecoder()
let data = jsonString.data(using: .utf8)!
let result = try decoder.decode(WeatherResponse.self, from: data)
print(result.weather[0].description)   // "晴"
\end{lstlisting}

铁律：\textbf{先看真实返回的 JSON，再定义字段}。很多“解析失败”的 bug 都是字段名写错、类型猜错（服务器发 “26.5” 字符串，你声明成 Double）——把原始 JSON 打印出来对照，十秒定位。

\booksection{五、URLSession + async/await：和服务器说话}

iOS 联网用系统自带的\textbf{URLSession}，配合 Swift 的\textbf{async/await}（异步函数），“等网络”就像“等快递”一样自然：

\begin{lstlisting}[language=Swift]
struct WeatherService {
    // 异步函数：async 标记"我要等一会儿"，throws 标记"可能失败"
    func fetchWeather(city: String) async throws -> WeatherResponse {
        let url = URL(string:
            "https://api.example.com/weather?city=\(city)")!
        let (data, _) = try await URLSession.shared.data(from: url)
        return try JSONDecoder().decode(WeatherResponse.self, from: data)
    }
}

// 调用处（在 .task 里）：
.task {
    do {
        let weather = try await service.fetchWeather(city: "北京")
        // 更新界面……
    } catch {
        print("出错了：\(error)")
    }
}
\end{lstlisting}

\textbf{避坑提醒：两条网络铁律}\par

(1)\textbf{联网必须 async/await，绝不阻塞主线程}——卡住界面是崩溃级体验。Swift 的结构化并发让这变得很自然，但别忘了\texttt{try await}和\texttt{do/catch}。(2)\textbf{ATS 默认禁止明文 HTTP}：iOS 要求 HTTPS。测试本地 http:// 服务时需要在 Info.plist 里临时开例外（NSAppTransportSecurity），但正式产品必须全走 HTTPS——审核也会盯这个。

\begin{tcbpractice}{来动手：做一个"每日一句"小应用}
\booksection{六、网络错误也要“分级”：别把一切报错都丢给用户}

断网、超时、404、服务器 500……不同的错，用户该看到的话不一样：

\begin{lstlisting}[language=Swift]
func friendlyMessage(_ error: Error) -> String {
    if let urlError = error as? URLError {
        switch urlError.code {
        case .notConnectedToInternet: return "网络断了，检查一下连接？"
        case .timedOut:               return "请求超时，稍后再试？"
        case .cannotFindHost:         return "服务器找不到了"
        default:                      return "网络出问题了：\(urlError.code)"
        }
    }
    if let d = error as? DecodingError {
        return "数据格式不对（可能是接口变了）"
    }
    return "出错了：\(error.localizedDescription)"
}
\end{lstlisting}

给用户看“人话”，给日志看“原话”（print 完整 error）。别把\texttt{error.localizedDescription}的英文天书直接甩给用户——产品感就藏在这些细节里。

找一个免费的名言 API（比如 hitokoto 之类），用 URLSession + async/await 拉取“每日一句”，显示在屏幕中央，加个“换一句”按钮。再把喜欢的句子用 SwiftData 存进“收藏夹”。这个练习把 UserDefaults/SwiftData/Codable/URLSession 全用了一遍——做完它，第 7、8 章就是小菜一碟。

\end{tcbpractice}

\begin{tcbsummary}{本章小结}
\begin{itemize}
  \item 口诀：小配置进 UserDefaults，成批数据进 SwiftData，外面的数据走 URLSession
  \item SwiftData 三板斧：@Model 声明表、modelContainer 挂容器、@Query 读 + context.insert 写
  \item Codable 自动"翻译" JSON：字段对得上，魔法就生效
  \item URLSession + async/await：联网像等快递；do/catch 处理失败
  \item 两条铁律：不阻塞主线程；正式产品必须 HTTPS（ATS）
\end{itemize}

\end{tcbsummary}
